\section{Context and goal}

Branch \code{darwin/gear\_radiation} implements a Fire-2-style
\citep{Hopkins2018} subgrid HII photoionization model for the GEAR
feedback scheme: each active star sorts nearby gas particles by
distance and spends a photon-rate budget on them, closest first, until
the budget is exhausted, tagging the affordable particles as
``ionized'' and holding their temperature at a floor value until the
next rebuild. Throughout this document, ``rebuild'' alone means a
star's own periodic HII candidate-search pass
(\code{GEARFeedback:HII\_rebuild\_time\_Myr}, typically sub-Myr); this
is a different, unrelated cadence from \swift's own tree rebuild
(\code{space\_rebuild}), which is always spelled out explicitly (e.g.\
``tree rebuild'') where it is meant. The model is validated against two
published, closely related reference algorithms:
\begin{itemize}
  \item \citet{Hu2017}, Appendix A: the four-source D-type expansion
    convergence test that our own ``spherical'' (\code{nside=0})
    mode reproduces directly;
  \item \citet{Smith2021}, Section 3.3.3--3.3.4 and Figs.~1--2: the
    same D-type test, plus an angular (HealPix) budget-splitting
    scheme designed to avoid ``mass-biasing'': a dense clump near a
    source absorbing a disproportionate share of the photon budget
    regardless of the small solid angle it subtends.

  Both papers are the ground truth this branch is being validated
  against; the example \code{examples/SubgridTests/\allowbreak
  SubgridRadiation/ClumpSmith2021/} reproduces their
  exact test geometry (four 19.2\,M$_\odot$ co-located sources,
  $n=100\,\rm cm^{-3}$ background, a $10^4\,\rm cm^{-3}$ clump
  20\,pc away).
\end{itemize}

\subsection{Scope of this document}
\label{sec:scope}

This is the \textbf{theory document}: it describes the current
algorithm (Section~\ref{sec:algorithm}) and how it is wired into
SWIFT's task graph (Section~\ref{sec:taskgraph}). It does not
track bug history, validation results, or open work. Those change
far more often than the algorithm itself, and outgrew what belongs
alongside shipped source. That material lives in the \textbf{GEAR
radiation logbook}, a separate document in the \code{swift-gear}
repository (\code{theory/radiation/}, branch
\code{darwin/radiation-logbook|main} once merged), covering:
\begin{itemize}
  \item a dated log of bugs found and fixed;
  \item the running validation campaign against Hu et al.\ (2017),
    Smith et al.\ (2021), STARBENCH, and the Grackle rate-coupled
    (T5) checks: tables, figures, and their interpretation;
  \item what remains undecided or undone (open items).
\end{itemize}
The two documents cross-reference each other by name and section
title rather than by automatic \code{\textbackslash ref}, since they
compile separately in separate repositories: a pointer like ``the
GEAR radiation logbook's Validation section'' is prose, not a live
link; if this document's section titles change, check the logbook for
stale pointers (and vice versa).

Update \code{01\_algorithm.tex} only when the code's physics actually
changes; update \code{02\_task\_graph.tex} only when the task-graph
implementation itself changes. Neither should accumulate validation
results or dated narrative; that belongs in the logbook.
